\subsection{Multiplication}

In this section, I want to explore how the add primitive can be used to build the meaning of multiplication. Multiplication is just repeated addition, so multiplication it is a great case study for how to build more meaning using narratives built from primitives. 

However, clearly I can establish multiplication on it own as a primitive because its essentially a primitive operator and there are clearly many optimizations, it is just a very simple case study that shows how meaning can be built on top of primitives.

So despite me establishing the meaning of multiply as a primitive, if I were to proceed with designing multiplication using repeated addition, I would do it as follows:

\begin{greyquote}
Multiply {multiplicand} and {multiplier}.
\end{greyquote}

\begin{greyquote}
Store 0 in result. Store 0 in counter. While the counter is less than multiplier, store result of add multiplicand and result in result then store result of add 1 and counter in counter.
\end{greyquote}

It is clear from this that to create the design that realizes this narrative, you need to establish how while will be automatically translated into a design structure. In this case, it is simply a conditional check that counter is less than the multiplier and apply the subsequent store operations. This can be automatically constructed as a design from the narrative and then an engine can reason about it. 

Note that I don't specify the invariants in the narrative, since the transformations are built from computable semantic primitives, the invariants are established through construction. 

Note that the conditional check in the while loop isn't a mechanical, it is inherently meaningful. It is reasoning about the state of the world and making a decision about which behavior to exhibit. If it fails, the meaning of that failure is unambiguous.

\subsection{Bill Total Example}

Using the add primitive, you can establish the meaning of behaviors. For example consider the behavior below in a restaurant payment app which is called add bill.

\begin{greyquote}
    Store 0 in bill-total. For each diner in the group, store the result of add diner-owes and bill-total in bill-total.
\end{greyquote}

Now in a larger narrative you can use add bill

\begin{greyquote}
    After the diners request the bill, store the result of add bill in total-owed.
\end{greyquote}

Since add bill is built from primitives that have unambiguous invariants, the correctness of the narrative can be established. For example if the each diners owes's type is not a number, the engine can reason that the add will fail.

